The Residual Triangle as Minimal Cell of an Equidistributed Orbit and the Propagation of Its Qualities and Quantities

1. The Minimal Cell

The residual near-equilateral triangle \(T\) is the smallest closed geometric object extracted from the infinite orbit generated by the practical step

\[
7\psi = \frac{\pi}{3} + \delta, \qquad \delta \approx 0.0051676144.
\]

Its three vertices are the even-index subsample

\[
V_k = \bigl(\cos(k\cdot 7\psi),\sin(k\cdot 7\psi)\bigr), \qquad k = 0,2,4.
\]

No smaller set of points drawn from the same orbit closes into a non-degenerate triangle that still carries a measurable six-fold memory. Hence \(T\) is the minimal cell.

2. Equidistribution under Further Iteration

The ratio \(7\psi/2\pi\) is irrational. By Weyl’s equidistribution theorem the sequence

\[
\theta_k = k\cdot 7\psi \bmod 2\pi, \qquad k = 0,1,2,\dots
\]

is therefore dense on the circle and equidistributed with respect to Lebesgue measure. Every arc of positive length is visited with asymptotic frequency proportional to its length. The residual \(\delta\) that prevents exact equilateral closure is precisely the arithmetic defect that forces this irrationality; the same defect therefore guarantees uniform distribution at every finer scale.

3. Propagative Quantities

All numerical content that appears at higher orders is already present, in concentrated form, inside the minimal cell:

| Propagative Quantity | Value | Origin inside \(T\) |
|----------------------|-------|---------------------|
| Residual seed \(\delta\) | \(0.005167614403\) | Angular defect of the generator |
| Thinnest side \(s_{\rm thin}\) | \(1.721623257385\) | Chord of even-index pair |
| Longer sides \(s_{\rm long}\) | \(1.737195272475\) | Remaining chords |
| Residual altitude \(h_{\rm res}\) | \(1.490969476641\) | Height relative to thinnest side |
| Shoelace area \(\operatorname{Area}(T)\) | \(1.2988990741\) | Exact Euclidean measure of the cell |
| Leftover area | \(1.8426935795\) | Unit-disk complement of \(T\) |
| Residual slice area \(A_{\rm slice}\) | \(0.6142311932\) | One-third of leftover |
| Dimensionless intensity \(\rho\) | \(0.472886\) | \(A_{\rm slice}/\operatorname{Area}(T)\) |
| Edge-curvature amplitude (bulge) | \(0.085119\) | \(c\cdot\rho\) with \(c=0.18\) |
| Inner-hexagon mean radius | \(0.577390\) | Intersection points of the two interlocking copies of \(T\) |

These quantities are invariant. They do not evolve; they are merely redistributed.

4. Propagative Qualities

Because the generator remains the same, every qualitative feature of the minimal cell is inherited by every later densification:

Six-fold memory.** The local angular spacing continues to oscillate about \(\pi/3\) with amplitude controlled by \(\delta\).  
Residual intensity.** The same \(\rho\) governs the curvature of every new edge and the metric of every new residual cell.  
Self-similarity.** Congruent copies of \(T\) (straight or curved) reappear at every radial scale and every order.  
Orientation.** The original counterclockwise-positive sense fixed by the Arcan(\(\tau\)) family is preserved under modular reduction.  
Phase-locking.** The entire densified structure remains synchronized with the discrete-tick lattice of the Cosmological Clock.  
Emergence and comoving transport.** The residual character that first appeared at Order-6 travels unchanged into Orders 10, 15, 30 and beyond.

5. Mechanism of Propagation

Further iteration simply appends more points of the identical orbit. At Order-10 the ten points still organise themselves into residual cells whose local geometry is a rigid image of \(T\). At Order-30 the thirty points fill the circle more finely, yet every triangular or hexagonal face continues to carry the same side-length ratio, the same residual altitude, the same intensity \(\rho\), and the same bulge amplitude. The hierarchy therefore contains no new metric invention; it contains only an increasing number of copies of the original minimal cell, packed at ever higher combinatorial density.

6. Dual Expression under Propagation

The two complementary operations remain active at every scale:

The winding method continues to read the altitude-to-terminal-ray gap and to return the intensity \(\rho\) as edge curvature.  
The radial-symmetry folding continues to close finite segments of the orbit into practical hexagons and to generate inner hexagonal measures by intersection.

Both operations propagate the same residual quantities and the same residual qualities outward through the nested shells.

7. Concluding Statement

The residual triangle \(T\) is the minimal cell of an orbit that becomes uniformly distributed under further iteration. All of its metric quantities—side lengths, altitude, area, residual intensity, bulge amplitude, inner-hexagon radius—and all of its qualitative features—six-fold memory, self-similarity, orientation, phase-locking, emergence—are propagative. They are fixed once and for all by the residual seed \(\delta\) and are thereafter only densified, never altered. The entire radial-fold hierarchy is the progressive, equidistributed unfolding of this single minimal cell.
